% =====================================================================
% Ensaio 1 — Seção 5: Dados, aquisição e descritivas
%
% Esta seção ESTÁ no alvo da qualificação. O roteiro pede "dose verificada,
% descritivas preliminares" — é isto. O que ela não traz é efeito estimado;
% esse material está na §6, e lá com o enquadramento que a pré-especificação
% impõe.
%
% ⚠️ Todo número aqui sai de artefato em data/processed/, produzido pelos
% scripts do pipeline, com vintage registrada. Nada é de memória.
% =====================================================================

\section{Dados, aquisição e descritivas}
\label{sec:dados}

\subsection{O que foi adquirido, e o que foi suspenso com razão declarada}
\label{sub:aquisicao}

A Tabela~\ref{tab:fontes} registra o estado de cada fonte. A distinção entre
\emph{não adquirida} e \emph{suspensa} não é burocrática: uma fonte suspensa foi
obtida, inspecionada e descartada por um defeito identificado, e o defeito é
informação sobre o desenho.

\begin{table}[htbp]
\centering
\caption{Fontes: estado em setembro de 2026}
\label{tab:fontes}
\small
\begin{tabular}{llp{7.2cm}}
\toprule
\textbf{Fonte} & \textbf{Estado} & \textbf{Observação} \\
\midrule
PAM/SIDRA        & adquirida & Tabelas 1612 e 1613; média 2015--2018 \\
SINASC           & adquirida & 2015--2022, via \texttt{pysus} \\
SIM/DOFET        & adquirida & Série nacional; o óbito fetal não está no DO \\
SINAN/IEXO       & adquirida & Substitui o SIH no canal de intoxicação \\
População (IBGE) & adquirida & SIDRA 6579; 2022 vem do Censo \\
FAO-GAEZ         & adquirida & 184 de 184 municípios \\
Bans municipais  & levantados & 17 tratados varridos; 1 achado \\
\midrule
SISAGUA          & \textbf{suspensa} & Quebra de registro laboratorial em 2020 \\
ANA (ottobacias) & \textbf{suspensa} & Bloqueada pela anterior, não por acesso \\
INMET/FUNCEME    & \textbf{suspensa} & Alísios: direção sem variação espacial \\
MapBiomas        & \textbf{descartada} & Banana cai em classe genérica \\
SEMACE (LAI)     & pendente & Única rota para o \(d=0\) operacional pré-ban \\
\bottomrule
\end{tabular}
\end{table}

\paragraph{As três suspensões, e por que cada uma é resultado.}
O SISAGUA foi obtido e tem \textbf{58.061} medições de agrotóxico no Ceará,
cobrindo os 184 municípios em todos os anos da janela --- a melhor cobertura de
qualquer fonte deste trabalho. Ele não entra mesmo assim, por três razões. A
primeira é que três das quatro moléculas que o Dossiê ABRASCO encontrou na
Chapada do Apodi --- procimidona, carbaril e fenitrotiona --- \textbf{não
constam da lista monitorada}: o que se mede bem é o glifosato, justamente o que
o Dossiê encontrou de menos. A segunda é que 95\% das amostras vêm do sistema de
distribuição, isto é, água já tratada, e apenas 0,4\% dos registros trazem
coordenada --- o desenho montante/jusante fica sem o insumo de que depende, e é
por isso que as ottobacias da ANA também ficam suspensas. A terceira é decisiva:
as detecções caem de 44 em 2019 para \textbf{zero} em 2020, 2021 e 2022, com
\emph{mais} amostras coletadas, enquanto o Brasil segue reportando cerca de
metade dos resultados em valor numérico. É mudança de prática laboratorial
cearense, iniciada no ano seguinte ao ban. Um DiD sobre essa variável concluiria
que a lei eliminou a totalidade das detecções, e o achado seria artefato.

O vento foi obtido do INMET e tem qualidade alta: quatorze estações automáticas
no Ceará, séries horárias com direção, velocidade e rajada integralmente
preenchidas. O teste que ele serviria --- separar deriva de confundidor pela
direção --- não funciona, e a razão é o contrário do que se esperaria: a
coerência circular \emph{entre} estações é de \textbf{0,96}. São os alísios, e
eles sopram sempre de leste. Com direção sem variação espacial, ``a favor do
vento'' é colinear com ``a oeste da fonte'', que é geografia --- e qualquer
confundidor com gradiente leste--oeste ficaria indistinguível de deriva.

\subsection{O painel}
\label{sub:painel}

O painel do Ensaio~1 é município \(\times\) ano-mês, balanceado: \textbf{184
municípios \(\times\) 96 meses = 17.664 células}, sem célula órfã. A
Tabela~\ref{tab:descritivas} traz as descritivas dos desfechos e das variáveis
de exposição.

\begin{table}[htbp]
\centering
\caption{Descritivas do painel, Ceará 2015--2022}
\label{tab:descritivas}
\small
\begin{tabular}{lr}
\toprule
\textbf{Grandeza} & \textbf{Valor} \\
\midrule
Nascidos vivos (pós-limpeza)            & 1.001.709 \\
Peso médio ao nascer                    & 3.208 g \\
Baixo peso (\(<\) 2500 g)               & 8,34\% \\
Prematuridade (\(<\) 37 semanas)        & 12,58\% \\
Óbitos fetais                           & 11.240 \\
Notificações de agrotóxico agrícola (SINAN) & 1.217 \\
População (Censo 2022)                  & 8.794.957 \\
\midrule
Municípios com dose \(> 0\)             & 169 \\
Municípios com dose \(= 0\)             & 15 \\
\bottomrule
\end{tabular}
\end{table}

\paragraph{O canal de intoxicação mudou de fonte, e a razão é aritmética.}
O desenho previa o SIH como contra-teste do canal-ar. Com a série completa de
oito anos, o SIH entrega \textbf{uma} internação por intoxicação acidental por
agrotóxico no estado inteiro --- e o placebo do canal, que contrasta acidental
com autoprovocada, morre junto, porque há um evento de cada lado. O SINAN/IEXO,
que é a notificação compulsória desenhada para exatamente isto, entrega
\textbf{1.217} notificações de agrotóxico de uso agrícola e 5.130 de
circunstância não intencional no mesmo período. Com um evento o canal não é
estimável; com 1.217 é. O SIH permanece como série secundária: ele mede
\emph{caso grave internado}, que é outro estimando, não um pior.

\paragraph{Duas limitações do SINAN que ficam declaradas.}
A circunstância --- o campo que sustenta o placebo --- está ignorada em
\textbf{32,7\%} das notificações, de modo que o contraste existe sobre dois
terços do dado. E o agente de uso em saúde pública, que seria a medida direta da
contaminação do grupo \(d = 0\) por controle vetorial aéreo (§2º do art. 28-B),
soma \textbf{89} notificações em oito anos. A flag da contaminação do zero
continua, portanto, sem fonte que a meça.

\subsection{A dose: o Gate 1 e a cultura-âncora}
\label{sub:dose-gate1}

A escolha da cultura-âncora era hipótese, não decisão, e o material de projeto
apontava para melão e algodão pela Chapada do Apodi. O dado a encerrou: o melão
tem área positiva em \textbf{10} municípios e o algodão em \textbf{28} ---
suporte insuficiente para a curva não-paramétrica, que exige valores de dose
espalhados e não empilhados. A banana sustenta: \textbf{169} municípios com área
positiva, Gini de \textbf{0,86} e 78\% da área no decil superior, que é a
combinação de dispersão alta com suporte largo que o \emph{sieve} pede. Ela
também é a cultura que casa com a química documentada --- procimidona é
fungicida de bananal, e aparece em 23 das 23 amostras do Dossiê.

\paragraph{O instrumento tem primeiro estágio.}
O índice FAO-GAEZ de aptidão agroclimática foi construído para os 184
municípios pela receita de \citet{reynier2025glyphosate}: diferença de
rendimento atingível entre alto e baixo insumo, percentil \emph{nacional} dessa
diferença --- calculado sobre os 5.570 municípios brasileiros, não sobre o Ceará
---, e máximo entre culturas tomado \emph{depois} do percentil. O decil superior
da banana tem aptidão média de \textbf{0,37} contra \textbf{0,26} do restante
(mediana 0,50 contra 0,20). Banana é plantada onde banana é apta, e a condição
de relevância se verifica por medida. \textbf{Isso não testa a exclusão}, que
não é testável.

\subsection{Quem de fato pulverizava por avião}
\label{sub:equipamento}

A dose é área plantada, e área plantada é \emph{proxy} do método. O Censo
Agropecuário permite olhar o método diretamente: ele registra o número de
estabelecimentos com uso de agrotóxicos \textbf{por tipo de equipamento de
aplicação}, com a aeronave como categoria própria, e desce ao município.

\begin{table}[htbp]
\centering
\caption{Equipamento de aplicação de agrotóxicos, Ceará (Censo Agropecuário 2006)}
\label{tab:equipamento}
\small
\begin{tabular}{lrr}
\toprule
\textbf{Equipamento} & \textbf{Estabelecimentos} & \textbf{\%} \\
\midrule
Pulverizador costal                & 105.624 & 95,75 \\
Tração mecânica e/ou animal        &     436 &  0,40 \\
\textbf{Por aeronave}              & \textbf{36} & \textbf{0,03} \\
\midrule
Total com uso de agrotóxico        & 110.312 & 100,00 \\
\bottomrule
\end{tabular}
\end{table}

\paragraph{A pulverização aérea é extraordinariamente concentrada.}
Trinta e seis estabelecimentos no estado inteiro, distribuídos por sete
municípios --- e dois deles respondem por 75\%: Limoeiro do Norte, com 18, e
Quixeré, com 9. São exatamente os dois municípios da Chapada do Apodi de que
trata o Dossiê ABRASCO.

\paragraph{Uma boa notícia para a identificação.}
\textbf{Nenhum município de dose zero tinha aeronave.} A flag da contaminação do
grupo de comparação --- a possibilidade de que municípios sem agricultura
pulverizada fossem tratados por dispersão aérea de controle vetorial --- não se
materializa nesta medida. É a verificação mais direta dessa ameaça que este
trabalho conseguiu produzir, e ela passa.

\paragraph{E uma má, que reinterpreta a §\ref{sec:diagnostico}.}
Dos dezessete municípios do grupo tratado --- o decil superior da banana entre
os 169 produtores ---, \textbf{quinze não tinham nenhum estabelecimento com
aplicação por aeronave}. O grupo concentra 27 dos 36 estabelecimentos, mas o faz
atribuindo dose alta a municípios onde não havia pulverização aérea a remover.

\begin{quote}
\textbf{Isso atenua o efeito estimado por construção, e é problema de definição
de tratamento --- não de poder estatístico.} Parte do que a
§\ref{sec:diagnostico} lê como ``o desenho não distingue o efeito do zero'' pode
ser, em alguma medida, ``o tratamento está medido no lugar errado''. As duas
leituras têm remédios opostos: a primeira pede mais unidades, a segunda pede
outra variável de dose.
\end{quote}

E o problema não se agrava com o da §\ref{sub:bans}: dos 27 estabelecimentos
do grupo tratado, 18 estão em Limoeiro do Norte, cuja lei municipal de 2009 foi
revogada em 2010. Os 27 --- Limoeiro do Norte e Quixeré --- são tratamento
genuíno em 2019, e a fração genuinamente tratada do grupo é de \textbf{dois
municípios em dezessete}.

\paragraph{Três ressalvas, e são sérias.}
O dado é de \textbf{2006}, treze anos antes do banimento, e a aviação agrícola
brasileira cresceu no período; o Censo de 2017 \textbf{não publicou} esse corte,
o que impede atualizá-lo. A medida conta \emph{estabelecimentos}, não área
tratada nem número de voos, e trinta e seis estabelecimentos podem cobrir área
considerável. E é declaração do informante sobre prática regulada, com
subdeclaração provável. Nada disto substitui a variável de dose: o que oferece é
uma medida direta daquilo que a dose aproxima, e a distância entre as duas é a
informação.

\subsection{Bans municipais anteriores: a ameaça medida}
\label{sub:bans}

O art.~29 da Lei Estadual nº 12.228/1993 autoriza municípios a legislar
supletivamente sobre agrotóxicos, e o material de projeto registrava que
Limoeiro do Norte teria proibido a pulverização aérea em 2009. Se procedesse,
haveria unidades já tratadas dentro do grupo de dose alta, e 2015--2018 deixaria
de ser período pré-tratamento para todas --- uma falha que o pipeline não
acusaria, porque o resultado sairia normalmente.

Procede. A \textbf{Lei Municipal nº 1.478, de 20 de novembro de 2009}, de
Limoeiro do Norte, dispõe sobre a proibição do uso de aeronaves nas
pulverizações de lavouras, nove anos antes da lei estadual. E o município está
no decil superior da banana --- sexto entre 169 ---, isto é, \emph{dentro} do
grupo tratado.

O levantamento dos 17 municípios do decil superior está completo: \textbf{um
confirmado}, dezesseis com acervo conferido e sem norma do tema, nenhum
inconclusivo. A única lei achada é a de Limoeiro do Norte, e o quadro não
cresce: a hipótese de que a Chapada do Apodi tivesse legislado em bloco não
procede --- Quixeré, com 21.289 leis no acervo digitalizado, e Russas, com
6.901, não têm norma anterior a 2019 sobre o tema.

\paragraph{E a lei foi revogada.} Ela durou seis meses. Em maio de 2010, um
mês depois do assassinato de José Maria Filho, o Zé Maria do Tomé, que a havia
defendido, a Câmara aprovou a \textbf{Lei Municipal nº 1.511, de 26 de maio de
2010}, que dispõe sobre a política ambiental do município e, num de seus
artigos, revoga a proibição\footnote{A ementa da Lei nº 1.511 não menciona a
revogação, que está no corpo da lei; por isso a busca por número e por termo no
acervo digitalizado da Câmara não a encontra. A ligação entre as duas leis vem
da imprensa: o \emph{Diário do Nordeste} noticiou, em abril de 2010, que o
projeto de política ambiental ``tem em um de seus artigos a revogação da lei
anterior, que proíbe a pulverização aérea''
(\url{https://diariodonordeste.verdesmares.com.br/editorias/regiao/comissao-vai-apurar-conflitos-na-chapada-do-apodi-1.102378}),
e o G1 (31/10/2024) identifica a lei votada em maio de 2010 como a nº 1.511
(\url{https://g1.globo.com/ce/ceara/noticia/2024/10/31/ze-maria-do-tome-o-ambientalista-cearense-assassinado-por-combater-o-uso-de-agrotoxicos-no-ceara.ghtml}).
A Comissão Pastoral da Terra (nota de 23/04/2014, reproduzida pela Terra de
Direitos:
\url{https://terradedireitos.org.br/noticias/noticias/quatro-anos-do-assassinato-de-ze-maria-uma-luta-contra-os-agrotoxicos-e-por-justica/14217})
e o MST (16/01/2019:
\url{https://mst.org.br/2019/01/16/proibicao-da-pulverizacao-aerea-de-agrotoxicos-no-ceara-direito-e-conquista-dos-povos-do-campo/})
datam a revogação de 20 de maio de 2010, provavelmente o dia da votação. O
texto integral da Lei nº 1.511 ainda não foi conferido.}. No pré-período de
2015--2018, portanto, Limoeiro do Norte não estava sob proibição: pulverizava
como os vizinhos, e a contaminação que esta subseção existe para medir
\textbf{não alcança a janela do estudo}. O levantamento vale pelo que mostra do
resto --- nenhum outro município do grupo tratado legislou --- e pelo que ensina
de método: lei achada não é lei vigente, e a varredura precisa procurar também a
revogação.

\paragraph{O que ``ausente'' significa aqui.}
Conferiu-se o acervo \emph{publicado} de cada câmara. Câmara que não digitalizou
2009 devolve vazio legitimamente, e dois dos dezesseis têm acervos pequenos
demais para serem completos. A tabela de ameaças (§\ref{sec:identificacao})
recebe isso como limitação declarada, não como zero.
